\documentclass[a4paper, 12pt]{report}

\usepackage[spanish]{babel}
\usepackage[utf8]{inputenc}
\usepackage[left=3cm, right=3cm, top=3cm, bottom=3cm]{geometry}
\usepackage{textcomp}
\usepackage{booktabs}
\usepackage{amssymb}
\usepackage{bussproofs}
\usepackage{fancyhdr}
\usepackage{graphicx}
\usepackage{amsmath}
\usepackage{enumitem}
\usepackage{ifsym}
\usepackage{float}


\usepackage{hyperref}
\hypersetup{
    colorlinks=true,
    linkcolor=blue,
    filecolor=magenta,
    urlcolor=cyan,
}

\pagestyle{fancy}
\lhead{Almeida, Figueroa, Ibarra \& Sánchez}
\chead{Tarea 6}
\rhead{\today}

\begin{document}
\begin{titlepage}
    \centering
    {\scshape\Huge Universidad Nacional Autónoma de México \par}
    \vspace{1.25cm}
    {\scshape\huge Fundamentos de Bases de Datos\par}
    \vspace{1.25cm}
    {\huge\bfseries Tarea 6:\\ Lenguaje de consulta SQL\par}
    \vspace{1.25cm}
    {\Large\textsc Almeida Rodríguez Jerónimo\par}
    \vspace{.1cm}
    {\large\texttt{418003815}\par}
    \vspace{0.25cm}
    {\Large\textsc Figueroa Sandoval Gerardo Emiliano\par}
    \vspace{.1cm}
    {\large\texttt{315241774}\par}
    \vspace{0.25cm}
    {\Large\textsc Ibarra Moreno Gisselle \par}
    \vspace{.1cm}
    {\large\texttt{315602193}\par}
    \vspace{0.25cm}
    {\Large\textsc Sánchez Lara Luis Enrique \par}
    \vspace{.1cm}
    {\large\texttt{418002856}\par}
    \vspace{1.5cm}
    \vfill
    \begin{figure}[hb!]
        \includegraphics[width=.3\textwidth]
            {../../logos/escudo_f-ciencias.png}\hfill
        \includegraphics[width=.3\textwidth]
            {../../logos/Escudo_UNAM.png}\hfill
    \end{figure}
\end{titlepage}

\subsection*{Inciso 5}
Las ventajas incluyen dar consistencia a los datos al asegurar que cierto valor
tiene que estar en la base de datos.\\
Cómo desventaja tenemos que al momento de poblar la base, si alguna de estas
restricciones no se cumplen entonces no van a poder ser ingresados los datos.
\end{document}
